% Milestone 3: energy efficiency, static analysis, model testing, data testing.
% Owner: @michudud04. Everything below is a draft assembled from the
% repository's own evidence so that its owner has a starting point, not finished
% text. The \tbd right after the heading says so in the PDF as well, and `make
% report-submit` counts it, so the draft cannot reach a submission unnoticed.
\subsection{Milestone 3: Model Building -- Quality Assurance}
\label{sec:m3}
\guidance{
  \textbf{Template:} Describe how the different software engineering practices and tools have been applied.\\
  The testing strategy for code, data and model is a typical M3 EDN entry.
}

\tbd{DRAFT, owner \texttt{@michudud04}: to be checked and rewritten in the owner's own words; the energy part waits for \#38.}

\subsubsection{Energy efficiency awareness}
\label{sec:m3-energy}
\guidance{
  \textbf{Template:} e.g.\ CodeCarbon. Reason about the energy efficiency of your ML component.\\
  \textbf{Rubric (5 pts):} How do our energy tracking practices help us reflect on the efficiency and sustainability of our choices?\\
  The demo wraps only the \texttt{fit} call in CodeCarbon's \texttt{EmissionsTracker} and logs the result to MLflow per run.
  Go beyond reporting a number: compare models or configurations by energy versus accuracy.
}

Energy is a requirement rather than a report item.
NFR-10 asks for three things: every training run records its emissions, training the chosen configuration takes at most 15 minutes on a laptop CPU without a hyperparameter search, and serving is reported as an average energy cost per answer.
Training is measured with CodeCarbon around the fit call and logged to the same MLflow run as the accuracy of that run, so the two are always read together.

Two decisions about how to measure are ours rather than the demo's.

\emph{The ladder is compared on energy, not only on error.}
Because the four variants of the experiment ladder differ in exactly one dimension at a time, the emissions figures answer a question the accuracy figures cannot: whether a rung earned its cost.
The extended feature set is the case this was built for.
It adds the four equipment lists as multi-hot columns, the history flags, four appearance columns and eight further technical columns, so it is strictly more expensive to build and to train than the basic set, and NFR-10's 15-minute budget is what it has to fit inside.
If it buys no accuracy, it is dropped, and the decision rests on a measurement of both axes rather than on the accuracy alone.

\emph{Serving energy is reported per answer and not per request.}
CodeCarbon's sampling granularity does not match single-digit-millisecond events, and a per-request tracker would eat into NFR-02's 200\,ms latency budget, so the serving figure is derived as an average over the NFR-03 load test instead of measured per call.
That is a deliberate loss of resolution, and it is written into the model card so that nobody reads the figure as something it is not.
The same reasoning already constrains the design elsewhere: the serving image carries no GPU or deep-learning dependency, \texttt{shap} is kept out of it because importing the module alone pulls in 189\,MB of \texttt{numba} and \texttt{llvmlite} (EDN-11), and multilingual text embeddings were left off the experiment ladder because they conflict with a CPU-only image.

The figures themselves do not exist yet:
\tbd{emissions and energy per training run for \texttt{b0}, \texttt{b1}, \texttt{lgbm-basic} and \texttt{lgbm-extended}, from CodeCarbon through MLflow (issue \#38)}
\tbd{the energy-versus-error comparison across the ladder, and what it decided about the extended feature set}
\tbd{the training time of the chosen configuration against NFR-10's 15-minute budget, and the average energy per answer from the NFR-03 load test}

\subsubsection{Static code analysis}
\label{sec:m3-static-analysis}
\guidance{
  \textbf{Template:} e.g.\ pylint.\\
  \textbf{Rubric (10 pts, shared with model and data testing):} To what degree does our linting contribute to the overall quality and consistency of the codebase?\\
  The lab lists Pynblint (notebook and repository QA), Pylint or flake8.
  We run ruff through pre-commit and CI; the demo additionally enables Pylint rules in ruff (\texttt{PL}), which our config does too.
}

Static analysis is ruff, as both linter and formatter, with the rule set \texttt{E, F, UP, B, SIM, I, PL}: pycodestyle errors, Pyflakes, pyupgrade, bugbear, simplify, import sorting and the Pylint-derived rules that cover the rubric's \enquote{Pylint or flake8}.
The whole configuration lives in \texttt{pyproject.toml}, and ruff also lints the code cells of every notebook.
On the commit this report was built from, \texttt{ruff format --check} reports 69 files already formatted and \texttt{ruff check} reports no findings.

Enabling the full rule set found two things.
The rules produced exactly five findings, all \texttt{PLR2004} on the same Cookiecutter placeholder line in the five CLI stubs; the line was removed and the rule stayed enabled.
More important was a gap in the configuration: its \texttt{include} key limited \texttt{make lint} to \texttt{recommenditos/}, but ruff does not apply \texttt{include} to paths passed explicitly, as pre-commit passes them, so CI linted \texttt{tests/} and \texttt{notebooks/} while the command every contributor runs did not.
The key is gone, and the property we have kept since is that the pre-commit hook, CI and \texttt{make lint} lint the same files, so a green \texttt{make lint} means a green CI job.

Three exclusions exist, each justified in the configuration itself.
\texttt{reports/analysis/} is excluded because those scripts are the evidence behind EDN entries and are kept exactly as they were run.
\texttt{PLR2004}, the magic-value rule, is ignored under \texttt{tests/}, because asserting against a literal expected value is what an assertion is.
\texttt{E402} is ignored in \texttt{recommenditos/tracking.py} alone, because mlflow reads \texttt{MLFLOW\_DISABLE\_AGENT\_HINT} at import time, so the variable has to be set before the import.

\paragraph{Notebooks: ruff, Pynblint and nbstripout (EDN-64, EDN-65).}
Our notebook is the exploratory analysis behind the dataset card: it recomputes the card's figures from the pinned source file through the pipeline's own functions, 174 checks in all, fails when one does not reproduce, and in being written corrected three of the card's statements.
Ruff checks the code in its cells, Pynblint its structure and narrative, and nbstripout removes outputs and execution counts on every commit, which also keeps rows of a dataset with personal data out of git.

Pynblint 0.1.6 \cite{quaranta2022pynblint} could not run as planned through a bare \texttt{uvx}: it pins \texttt{typer<0.13} but not click, and that typer crashes on click 8.2 and newer.
It runs instead from a hash-locked environment of its own in \texttt{tools/pynblint-env/}, apart from the project's lock.
It also exits 0 whatever it finds, so \texttt{tools/notebook\_lint.py} is the gate: it runs Pynblint in a sanitised environment, because Pynblint reads its settings from unprefixed environment variables, then reads its JSON report and fails on any finding.
CI runs the same \texttt{make notebook-lint}.

The policy is a table of all 22 rules in the documentation, which the gate reads and which must classify exactly the rules the installed version has.
Thirteen rules are enforced at their default thresholds.
Three read execution counts, which nbstripout removes, so on a committed notebook one would always fire, one never, and one only repeat \texttt{empty-cells}; they are excluded, and \texttt{make notebook-run} shows instead that the notebook runs top to bottom, outside CI because it needs the 548\,MB source file.
The five repository rules are not run, because each checks something enforced more strongly here, such as a \texttt{.coverage} file where CI publishes the coverage of every run, and \texttt{notebook-name-too-long} is disabled in Pynblint itself.
Two rules of our own add what Pynblint does not check: every notebook lives in \texttt{notebooks/} under the Cookiecutter name, and none carries the \texttt{keep\_output} flags that would make nbstripout commit its outputs.
A rule that is wrong for us is excluded in the table, with its reason, in a reviewed pull request; there is no warning tier and no waiver.

On its first run Pynblint found five code cells over the 30-line limit, which were restructured rather than the limit raised; on the stripped copy that is committed, only \texttt{non-executed-notebook} fired, which is the measurement behind excluding the execution rules.
On the commit this report was built from it reports no findings.

\subsubsection{Model testing}
\label{sec:m3-model-testing}
\guidance{
  \textbf{Template:} e.g.\ pytest.\\
  \textbf{Rubric:} To what extent do our tests give confidence that the code and model behave as expected under different conditions?\\
  The demo tests a metric threshold on the test set and uses parametrized invariance tests.
  Optional extension listed by the demo repo: non-functional requirements (explainability with SHAP, fairness with AIF360, trustworthiness with TrustML).
}

The suite is Pytest with coverage measured on every run.
On the commit this report was built from it is 674 tests, 670 passed and 4 skipped where they cannot apply, in about three minutes, at 98\,\% statement and branch coverage of \texttt{recommenditos/}.
Coverage measures branches rather than only statements, so a half-tested conditional does not count as covered, and nothing is excluded from measurement except the \texttt{if \_\_name\_\_ == "\_\_main\_\_"} entry points, which Pytest never executes.
One module sits at 0\,\%, \texttt{plots.py}, a Cookiecutter template stub that nothing imports yet; it is counted rather than excluded, so the number does not flatter itself.
There is deliberately no \texttt{--cov-fail-under}: NFR-07 asks for 80\,\% on a delivery commit, which is when the number is read and acted on, and a per-pull-request gate would block merges while the suite is still being built.

Three things go beyond asserting a threshold.

\emph{Tests carry the requirement they verify.}
A test declares its requirements with the marker \texttt{@pytest.\allowbreak mark.\allowbreak req("NFR-08")}, and the marker is what CI turns into a requirement-to-test matrix, so traceability is generated from the tests rather than maintained as a table that goes stale (NFR-07).
Five requirements carry markers today, NFR-01, NFR-06, NFR-07, NFR-08 and NFR-14, across 63 tests; the matrix itself is built and published by CI on every run, and its gate is scoped to the milestone that owes each requirement its evidence (EDN-44).
What a marker earns depends on what the specification says the evidence is: on an \textbf{[automated]} entry a test is that evidence, while on a \textbf{[manual]} entry the promised evidence is the drill the entry names, so a marker there records that a test touches the requirement and does not stand in for the drill (EDN-45).
That is why the matrix reports NFR-06 as named by a test rather than verified by one: the download stage's pin and MD5 tests are evidence about reproducibility, and the evidence NFR-06 promises is a \texttt{dvc repro} re-run on a clean clone.
No status in the matrix claims that a person judged the evidence sufficient, which is why NFR-07 keeps that review before each delivery.

\emph{The suite never runs the pipeline through DVC.}
It calls the same stage functions in the same order against the generated synthetic fixture in a temporary directory, so it is hermetic and needs no credentials, no 548.6\,MB download and no network (EDN-33).
Every edge case the pipeline rules exist for is guaranteed present in that fixture whatever the row count, so a test asserts against it rather than hoping the data contains it: a listing registered after the snapshot, a duplicate on the deduplication key, a Spanish listing, a make far below the support threshold, prices outside the training range at both ends, a new car, a pre-registered car, a transporter, a missing registration date, a private seller, and a used car whose flag says otherwise.
The fixture also keeps the real make skew and gives each seller several listings, because a seller-grouped split on a flat distribution would be indistinguishable from a random one.

\emph{One test exists because a test was wrong.}
The fixture's first implementation built each edge-case row as a copy of one body row, which made most of them duplicates of it on the deduplication key and gave them all that row's country, so preprocessing would have deleted several of the cases the fixture exists to provide and diverted the rest into the Spanish holdout.
The first round of tests could not see it, because they asserted that each case was present in the raw frame, which was true.
The fix carries a test that asserts each case \emph{survives} the preprocessing rules instead, and the lesson, that a check has to sit where the value is consumed, is now in the working agreements.
The same reasoning produced a second pair of tests around the quality gate: one asserts that the stub model does not pass it, and a second asserts that the gate still discriminates, because a gate that says no to everything would satisfy the first test.

The model tests that came with the train and evaluate stages test the gate rather than the model's numbers, which \texttt{evaluate} measures on the real test split (\cref{sec:goal}): 21 tests carry NFR-01 and drive each criterion through its pass, its fail and its not-measured path with constructed models and records, SC-04's to SC-06's bounds at their edges, and others hold \texttt{predict\_eur} to its contract, such as one row predicting what the batch predicts for it and every prediction staying inside the training price range.
Invariance tests under small input perturbations, and the check that the booster's own SHAP export matches \texttt{shap.TreeExplainer}, which is what makes it safe to keep \texttt{shap} out of the serving image (EDN-11), are still to be written.

\subsubsection{Data testing}
\label{sec:m3-data-testing}
\guidance{
  \textbf{Template:} e.g.\ Great Expectations, Deepchecks.\\
  \textbf{Rubric:} How effectively does our data validation strategy anticipate and prevent potential issues in the pipeline?\\
  The demo validates raw and clean data with Great Expectations as a DVC stage and publishes Data Docs.
}

Data validation is split into two layers on purpose, because structure and value answer different questions (EDN-31).
A structural break is a bug in our code and must fail the stage that caused it; a value break is a change in the data and belongs in an expectation suite with a tolerance.

The structural layer is \texttt{recommenditos/schema.py}: three frozen contracts fixing the column names, dtypes and nullability of the raw, interim and processed frames, with a \texttt{validate} that reports every problem at once with the offending row positions, and a \texttt{conform} that selects and casts before each write.
It is hand-written rather than taken from a validation library, and that was decided on measurements rather than on documentation: a Parquet round trip silently turns an \texttt{object} string column into \texttt{str}, so a contract declaring \texttt{object} passes upstream and fails downstream, and Pandera does not check the datetime unit at all, so a \texttt{datetime64[ns]} declaration there would be documentation rather than a guarantee.
Because \texttt{conform} selects only the contract's columns, the seven personal-data columns NFR-08 forbids cannot reach a processed artefact even if a stage forgets to drop them, and two tests assert that they do not.

The value layer is Great Expectations, in two DVC stages that fix the demo's two warts (EDN-68).
\texttt{configure\_gx} builds the context, and its output \texttt{gx/} is a dependency of \texttt{validate-data}, so validation cannot run against a missing or stale store; it is cached rather than committed, because Great Expectations writes fresh UUIDs on every save.
The suites are generated, their schema checks from \texttt{schema.py} and every bound from \texttt{params.yaml}; a row minimum and fill rates on the bounded columns keep any rule from passing vacuously, and each rule sits on the frame where it can hold.
\texttt{registration\_date} at or before the reference date is hard on the cleaned frame and tolerated on the raw one up to \texttt{mostly=0.99}, so a future scrape carrying many such rows is flagged rather than silently cleaned (EDN-22); the training price range is checked on the cleaned frame only, since the published file starts at 1\,EUR; and the cleaned frame's column list is NFR-08's check that no personal-data column survives.
Two planned rules were dropped on evidence: an \texttt{is\_used}/\texttt{offer\_type} agreement fails 18,108 rows of the published file by design (EDN-23), and the supported-make list only exists after this stage (EDN-48).
Probed per expectation on our pandas-3 frames, the tool works, but compares only a column's scalar type, so it cannot tell nullable, extension or \texttt{object} dtypes from their counterparts, which \texttt{schema.py} owns, nor bound a text date, so the raw suite sees that column parsed.

On the real snapshot, the raw suite's 103 expectations find the 164 registrations after the reference date, 0.14\,\% of the dated rows and inside its tolerance, and three odometer readings above 1,000,000\,km, the highest 2,570,000\,km.
All 57 of the cleaned suite pass, the mileage rule included (EDN-72): the three readings survive preprocessing, and the range is a check that tolerates 0.1\,\% rather than a filter.
A failed expectation fails the stage (\cref{sec:m2-data-versioning}) after writing the Data Docs, which \texttt{dvc pull} restores, and \texttt{reports/\allowbreak data-validation/\allowbreak summary.json}, git-tracked and free of timestamps, so its diff in a pull request is exactly what the data did.
